\documentclass{article}
\input{boilerplate}

\begin{document}


\newcommand{\lectureTitle}{YOUR TITLE HERE}
\newcommand{\lectureDate}{Due: May 11th, 10:30 PM}

\textsc{COS435 / ECE433: Introduction to RL}
\vspace{1em} \hfill \lectureDate

\maketitle

\paragraph{Authors.} Jane Doe (\texttt{jane.doe@princeton.edu}), John Smith (\texttt{john.smith@princeton.edu})

\paragraph{Link to GitHub Repo.} [Your link here. Make sure it is public!]
\section*{Instructions}

This is a rough template guide for your final project report. If you do not have these exact sections in your final report, or do not follow this formatting, do not worry -- prioritize the contents over the exact section breakdown here.

A few things to keep in mind:
\begin{itemize}
    \item {\bf The main body of your final report must be $\leq 8$ pages.} If your current report is longer than $8$ pages, please push finer details to an Appendix or cut down on the current text. You should {\it probably} aim to have at least 4 pages of content, but again, we won't automatically take off points if the length is too short.
    \item {\bf There is no minimum length for the final report.} Please make sure to cover the rubric items in the final project guidelines.
    \item {\bf The minimum allowed font size if 10 pt. as it is in this template for readability.} There is no restriction on the specific font you choose. Because there is no minimum length, you can choose to make your font-size 12 pt., but please keep everything reasonable-looking (i.e. no 16 pt. situation).
    \item {\bf Signpost.} Please do some signposting, even if you don't use the section breakdown here -- this is to ensure that we don't miss anything within your report. It is okay to even, i.e., include headings for your paragraphs:
    \\\\{\bf What is the problem? } Insert text here. \\

    but we encourage you to write this report like a realistic NeurIPS, ICLR, ICML, etc. paper where content is wrapped into the main sections. Again, we will evaluate the reports on {\bf content} over the exact formatting and structure.
\end{itemize}

\section{Introduction}

Introduce the problem. Please cover the following points:
\begin{itemize}
    \item Motivation: What is the problem?
    \item Motivation: Why is it interesting and important?
    \item Motivation/relationship with prior work: Why is it hard? (E.g., why do naive approaches fail?)
    \item Relationship with prior work: Why hasn't it been solved before? (Or, what's wrong with previous proposed solutions? How does mine differ?)
    \item Methods and results: What are the key components of my approach and results? Also include any specific limitations.
\end{itemize}

\section{Methods}

Write about your methods.

\section{Results}

Write about your results. Please include at least one new figure that are based on your own results.

\section{Discussion and Limitations}
Please include a discussion of the limitations of your work, the prior paper, and/or method.

{\footnotesize
% \paragraph{Acknowledgements.} We thank XXX.

\bibliographystyle{apalike}
\bibliography{references}
}

\section{Appendix (Optional)}

Include any details that don't fit within the main body here.

\end{document}